\begin{exercise}[Gauge-theory rules before gauge fixing]{I-CH08-003}
For canonically normalized fields, use
the mostly-plus metric and Clifford convention
$\eta=\operatorname{diag}(-,+,\ldots,+)$,
$\{\gamma^\mu,\gamma^\nu\}=-2\eta^{\mu\nu}$, and
$\slashed p=\gamma^\mu p_\mu$.
\[
 \mathcal L=-\frac14F_{\mu\nu}{}^aF^{a\mu\nu}
 -(D_\mu\phi)^\dagger D^\mu\phi-V(\phi)
 +\bar\psi(-\ii\gamma^\mu D_\mu)\psi-\mathcal L_{\rm Yuk},
 \qquad D_\mu=\partial_\mu-\ii gA_\mu{}^aT^a.
\]
Expand the action about a stationary vacuum, allowing it to break the gauge
group, and derive the unfixed quadratic kernels and the gauge, scalar, and
fermion interaction vertices.
It is enough to express the vertices from $V$ and $\mathcal L_{\rm Yuk}$ as
field derivatives of those functions.  With all momenta incoming, write the
three- and four-gauge-boson vertices and the one- and two-gauge-boson scalar
vertices explicitly.  Identify every gauge-orbit zero direction of the
quadratic action and show directly why the naive gauge-field propagator does
not exist before gauge fixing.
\end{exercise}
